The present work is set in two parts. Firstly, an analysis of different platforms for capturing aerial images of native forest in Misiones is carried on, for monitoring and preservation purposes. Comparisons of costs and features among different satellites and unmanned aircraft systems (UAS) are made, for different areas of native forest, either of public or private domain, from few hectares to huge extensions of thousands of hectares. In the second part, different ways of image processing are analyzed. Specifically three methods are revised, one of them is shadow detection by means of homomorphic filtering, another is segmentation by means of morphologic processing, and lastly the shadow detection by means of invariant index. For the three methods pros and cons and the viability for their application to the monitoring of native forest reserves in Misiones are analyzed. \\
\begin{keywords}
    Aerial image, Illumination Invariant Shadow Ratio, Image processing, Remote sensing
\end{keywords}    %Palabras Claves en Español (tres a seis palabras claves separadas por comas. La primera letra de cada palabra clave debe empezar con mayúscula) 

%Aerial image, Illumination Invariant Shadow Ratio, Image processing, Remote sensing